\section{Execution}
\label{chap:execution}
In this chapter, a description of how the testing of Hover was designed and conducted will be presented.

\subsection{Experiment Design}
The experiment was designed mainly to answer the research questions defined in Chapter \ref{chap:rqm}. Participants were picked using a non-probabilistic approach where we used a combination of the purposive and convenience sampling technique. Our priority was to get concise feedback from users with a behavioral pattern that could benefit from using Hover. Additionally, we wanted to pick groups of friends to test the social aspect of the game in a realistic manner. Thus we prioritized gathering participants we knew to ensure this. Triangulation with multiple types of data was used to improve the validity and reliability. Specifically, four different data types were gathered and analyzed: \textit{Questionnaire}, \textit{semi-structured interviews}, \textit{participant observations}, and \textit{application usage data}. A detailed description of the content and purpose of collecting these data types is further discussed in Chapter \ref{chap:data_collection}.

\subsection{Conducting the Experiment}
The experiment took place in the two weeks between the 12th and 25th of April 2021. We wanted to test the effects Hover may have on people’s habits, and consequently, we wanted a test period that was as long as possible. Due to limited time in the Master Thesis project, we decided that two weeks would be sufficient. 

Before the experiment started, the participants were asked to watch a quick introduction to Hover video. The purpose of this was to ensure users would begin on the same level of knowledge and illustrate all the game features. We initially planned to give the introduction physically, but because of the ongoing COVID-19 situation, we decided it was best to adapt the presentation to an online video instead.

To distribute information to the participants efficiently, we had created a \href{https://hover-game.web.app/testing}{website} to gather everything. The page included a what to do guide, links to the survey and introduction video, and other essential information. The application was distributed to iOS users via \href{https://apps.apple.com/no/app/hover/id1556425356}{App Store} and to Android users via \href{https://play.google.com/store/apps/details?id=com.hoverapp.hover}{Google Play}.

When the two weeks of testing came to an end, all participants were asked to complete a survey with questions based on our research questions. The survey took ten to fifteen minutes to answer and included an introduction question that should identify users that had not been using the application. These answers were discarded.

As gratitude for everyone taking their time to participate in this study, five surprise gifts were given to randomly picked participants. The gifts helped us ensure that all people participating in the survey would take part in the whole period and complete the questionnaire at the end. We choose to pick random winners not to stimulate the user’s interest in Hover.

\subsection{Summary}
In this chapter, a description of the experiment and how it was designed and conducted was given. Next, Chapter \ref{chap:data_collection}, will look at how and describe what types of data we collected in more detail.
